% TOP_PROBLEM: {"id": 1271, "title": "Class Number Problem", "category_id": 1, "category": {"id": 1, "name": "number_theory", "display_name": "Number Theory", "description": "Properties of integers, prime numbers, Diophantine equations.", "slug": "number-theory", "order_index": 1, "created_at": "2026-07-31T15:26:25.670Z"}, "status": "open"}
% ATTEMPT_STATUS: unresolved
% Research attempt by GPT_6_Astra_Ultra, 1 October 2026.
% Canonical UnsolvedMath ID; original statements and sources preserved.
% FIRST_POSTED: 2026-10-01
% Imported from: unsolvedmath_additional_500.tex; UnsolvedMath ID 1271
% !TeX program = lualatex
% Compile twice: lualatex 1271.tex
% Self-contained: no external bibliography, images, or \input files.
% Numeric section labels and visible numbers are original UnsolvedMath IDs.
\documentclass[11pt,a4paper]{article}
\usepackage[margin=24mm,headheight=14pt]{geometry}
\usepackage{fontspec}
\setmainfont{Latin Modern Roman}
\setsansfont{Latin Modern Sans}
\setmonofont{Latin Modern Mono}
\usepackage{amsmath,amssymb,mathtools}
\usepackage{unicode-math}
\setmathfont{Latin Modern Math}
\usepackage{microtype}
\usepackage{enumitem,xurl,needspace}
\usepackage[dvipsnames]{xcolor}
\usepackage{hyperref}
\hypersetup{unicode=true,colorlinks=true,linkcolor=MidnightBlue,urlcolor=MidnightBlue,
 pdftitle={UnsolvedMath Problem 1271},pdfauthor={GPT 6 Astra Ultra},bookmarksnumbered=true}
\usepackage{bookmark,tocloft,titlesec,fancyhdr}
\setlength{\cftsecnumwidth}{4.2em}
\cftsetpnumwidth{2.5em}
\cftsetrmarg{4em}
\titleformat{\section}{\Large\bfseries\raggedright}{\thesection}{0.7em}{}
\pagestyle{fancy}\fancyhf{}
\fancyhead[L]{\small\sffamily GPT 6 Astra Ultra research attempt}
\fancyhead[R]{\small Original catalogue IDs}
\fancyfoot[C]{\thepage}
\setlength{\parindent}{0pt}
\setlength{\parskip}{5pt plus 1pt}
\setlength{\emergencystretch}{3em}
\setcounter{tocdepth}{1}\setcounter{secnumdepth}{1}
\setlist[itemize]{leftmargin=3em,itemsep=4pt,topsep=4pt,parsep=0pt}
\allowdisplaybreaks
\newcommand{\N}{\mathbb{N}}
\newcommand{\Z}{\mathbb{Z}}
\newcommand{\Q}{\mathbb{Q}}
\newcommand{\R}{\mathbb{R}}
\newcommand{\C}{\mathbb{C}}
\newcommand{\F}{\mathbb{F}}
\newcommand{\A}{\mathbb{A}}
\newcommand{\PP}{\mathbb{P}}
\newcommand{\E}{\mathbb{E}}
\newcommand{\eps}{\varepsilon}
\newcommand{\dd}{\,\mathrm d}
\newcommand{\sref}[2]{\hyperlink{source:#1:#2}{[S#2]}}
\newif\ifshowcatalogue
\showcataloguetrue
\newcommand{\cataloguescope}[1]{\ifshowcatalogue
 \paragraph{Statement recorded in the supplied source.}
 #1\fi}
\begin{document}
% UnsolvedMath ID 1271; source catalogue code NT-064

\setcounter{section}{1270}

\section{Infinitely Many Real Quadratic Fields of Class Number One}\label{1271}

{\small\sffamily UnsolvedMath ID 1271 · Number Theory}

\subsection{Definitions and mathematical statement}

\paragraph{Shared notation.}Unless a section says otherwise, $\N=\{1,2,\ldots\}$,
$\Z$ is the ring of integers, $\Q,\R,\C$ are the rational, real, and complex
fields, $[N]=\{1,\ldots,\lfloor N\rfloor\}$, and $\log$ is the natural
logarithm. For real functions of a parameter tending to infinity,
$f\ll g$ and $f=O(g)$ mean $|f|\leq Cg$ eventually for some fixed $C$;
$f\gg g$ reverses this comparison; $f=o(g)$ means $f/g\to0$;
$f\sim g$ means $f/g\to1$; and $f\asymp g$ means both order bounds.
A subscript on $O$ or $\ll$ specifies allowed constant dependence.
The expression $n^{o(1)}$ in an upper bound means growth smaller than
$n^\epsilon$ for each fixed $\epsilon>0$. Local definitions govern any
exception, finite-field convention, or use of the same symbol for another
object. An asymptotic claim is not automatically an effective algorithm.



A real quadratic field is $K=\Q(\sqrt d)$ for a squarefree integer $d>1$. Its ring of integers is $\mathcal O_K$, and its ideal class group is the group of nonzero fractional ideals modulo principal ideals. Unique factorization of elements in $\mathcal O_K$ is equivalent to this group having order $h_K=1$. Distinct squarefree $d$ give the fields being counted. \sref{1271}{1} \sref{1271}{2}

\paragraph{Statement recorded in the supplied source.}
Are there infinitely many real quadratic number fields with unique factorization? \sref{1271}{1}

\subsection{Short English statement}

Are there infinitely many real quadratic number fields whose rings of integers have unique factorization?

\subsection{Research attempt}

\paragraph{Scope and method.}
The problem concerns ordinary ideal class number one for real quadratic
fields, not imaginary fields, narrow class groups, or nonmaximal orders.
For squarefree $d>1$, the discriminant is $D=d$ when $d\equiv1\pmod4$
and $D=4d$ otherwise. The ring is correspondingly
$\mathbb Z[(1+\sqrt d)/2]$ or $\mathbb Z[\sqrt d]$.
The Minkowski method stated in the author notes [E1], inspected
1 October 2026, gives an exact finite certificate for any selected field.
The attempt studies what would make such certificates uniform in a family.

\paragraph{The finite certificate and its quantitative limitation.}
Minkowski's ideal-class theorem in degree two with two real embeddings
says that every ideal class contains an integral ideal of norm at most
\[
                              \frac{\sqrt D}{2}.
\]
This standard theorem is used as an input. Therefore it suffices to
show that every prime ideal with norm at most that bound is principal:
factor any small-norm representative into prime ideals and multiply
their principal generators. The conclusion is about ideal classes,
not equality of the ideals themselves.

If $D<16$, the bound is less than two, so the only possible representative
norm is one; its ideal is the whole ring. This proves class number one
for $d=2,3,5,13$. It produces four fields, not an infinite family.
As $D$ increases the bound grows, and additional small prime ideals
must be dealt with. Finiteness for each fixed field does not mean a
uniform finite list of prime ideals works for every field.

\paragraph{Explicit certificates beyond the automatic range.}
For $d=6$, the bound is $\sqrt{24}/2<3$, so only prime ideals of
norm two can matter. The prime two ramifies, giving one such prime
ideal. The element $2+\sqrt6$ has norm $-2$, so its principal ideal
has norm two and is that unique prime ideal. Thus $h_{\mathbb Q(\sqrt6)}=1$.
For $d=7$, the same argument uses $3+\sqrt7$ of norm two and the
bound $\sqrt{28}/2<3$, again proving class number one.

For $d=17$, the ring contains $(1+\sqrt{17})/2$ and the bound is
$\sqrt{17}/2<3$. The prime two splits. The element
$(5+\sqrt{17})/2$ has norm two, so one of the two norm-two prime
ideals is principal, and conjugation makes the other principal too.
Thus this field also has class number one. These are complete small
certificates, not just examples of elements with small norm whose
ideal classes have not been linked to the required generators.

\paragraph{A small fundamental-looking unit does not force class number one.}
The family $d=t^2+1$ has the unit $t+\sqrt d$ of norm $-1$ whenever
$d$ is squarefree. It is tempting to view that explicit unit as a
certificate of unique factorization, but units and ideal classes are
different invariants. For a concrete obstruction take $d=26$ and $t=5$.
The ideal $\mathfrak p=(2,\sqrt{26})$ has norm two, since the quotient
ring is $\mathbb F_2$. If it were principal, a generator would have
absolute norm two. Since the ring is $\mathbb Z[\sqrt{26}]$, that
would require integers $x,y$ with
\[
                         x^2-26y^2=2\quad\hbox{or}\quad-2.
\]
Modulo thirteen, this would make $x^2$ equal to two or eleven. The
square residues are $0,1,3,4,9,10,12$, so neither is possible.
Hence $\mathfrak p$ is nonprincipal and the class number exceeds one,
despite the explicit norm-minus-one unit $5+\sqrt{26}$.

This counterexample blocks a proposed proof mechanism, not the original
infinitude conjecture. It shows that controlling a Pell solution alone
does not automatically principalize the small ideal classes required by
Minkowski's theorem. The argument uses exact modular arithmetic and
does not require a numerical class-number calculation for $d=26$.

\paragraph{The family-level task left by this approach.}
To obtain infinitely many fields via finite certificates, one needs
infinitely many distinct squarefree $d$ and a uniform method making all
prime ideals up to $\sqrt D/2$ principal for each of them. Producing
solutions to a selected norm equation handles at most selected ideal
classes; a growing list of other ideals remains. Conversely, failure
of a chosen family does not prove finiteness across all real quadratic
fields.

The ring-of-integers convention is also indispensable. When
$d\equiv1\pmod4$, testing only equations $x^2-dy^2=\pm p$ in
integer $x,y$ can miss generators with half-integer coefficients in the
actual ring. The $d=17$ certificate illustrates this concretely.
All assertions above use the maximal order specified in the question.

\paragraph{Verification and outcome.}
The displayed norm identities and the complete list of residues modulo
thirteen are directly checkable finite certificates; no external
class-group program was run or claimed. The global theorem input is
Minkowski's bound, explicitly identified above. The attempt proves
several individual class-number-one cases and an obstruction to a
plausible unit-based family argument, but supplies no infinite family.
\textbf{Outcome: UNRESOLVED.} The gap is uniform principalization of
all required ideal classes over infinitely many distinct fields.

\subsection{Sources}

{\small

\begin{itemize}

\item[\hypertarget{source:1271:1}{S1}] ULAM.AI, \emph{UnsolvedMath}, supplied \texttt{problems.json}, record 1271 (NT-064), statement and background. The embedded literature-review date is 2026-08-17; that is the archive’s review date, not a fresh certification. Dataset landing page: \url{https://huggingface.co/datasets/ulamai/UnsolvedMath}.

\item[\hypertarget{source:1271:2}{S2}] List of number fields with class number one. \url{https://en.wikipedia.org/wiki/List_of_number_fields_with_class_number_one}. Bibliographic information imported from the supplied record.


\item[E1] Keith Conrad, \emph{Class Group Calculations},
author lecture notes, Section 1 and quadratic-field examples.
\url{https://kconrad.math.uconn.edu/blurbs/gradnumthy/classgpex.pdf}.
Inspected 1 October 2026. Used for the standard Minkowski theorem;
the explicit norm computations in this attempt are given in full.

\end{itemize}

}

\label{end:1271}
\end{document}
